\documentclass[11pt]{article}

\usepackage[margin=1in]{geometry}
\usepackage{amsmath, amssymb}
\usepackage{listings}
\usepackage{xcolor}
\usepackage{hyperref}
\usepackage{enumitem}
\usepackage{titlesec}
\usepackage{parskip}
\usepackage{booktabs}

% Code listing style
\lstset{
    basicstyle=\ttfamily\small,
    keywordstyle=\color{blue}\bfseries,
    commentstyle=\color{gray}\itshape,
    stringstyle=\color{red},
    frame=single,
    breaklines=true,
    showstringspaces=false,
    tabsize=4
}

\title{\textbf{CS3210: Operating Systems} \\ Lecture 14 --- User and Kernel Threads}
\author{John Kim \\ Georgia Institute of Technology}
\date{March 2, 2026}

\begin{document}

\maketitle
\tableofcontents
\newpage

% ============================================================
\section{A Note on Weak Ordering}
% ============================================================

\subsection{Weak Order Semantics}

\textbf{Weak ordering} (also called \emph{relaxed consistency}) describes a property of modern hardware and compilers: they are permitted to reorder reads and writes as long as the result remains correct \emph{from the perspective of the executing thread itself}. The important caveat is that weak ordering does \emph{not} concern itself with correctness across concurrently executing threads.

In practice this means:
\begin{itemize}
    \item Stores performed by one thread can appear to arrive out of order as observed by another thread.
    \item A thread might see another thread's writes to different variables at different times, depending on caching and pipeline effects.
    \item Without explicit synchronization, almost no ordering guarantees exist between threads.
\end{itemize}

\subsection{Lock Acquire/Release as a Partial-Order Mechanism}

The primary tool for imposing order is the \textbf{lock acquire/release} pair:
\begin{itemize}
    \item \textbf{Acquire} $\rightarrow$ all subsequent reads and writes in the current thread are guaranteed to happen \emph{after} the acquire, as seen by any other thread that later acquires the same lock.
    \item \textbf{Release} $\rightarrow$ all preceding reads and writes in the current thread are guaranteed to happen \emph{before} the release, again as seen by any thread that later acquires the lock.
\end{itemize}

This gives us a \emph{partial} ordering of memory operations, not a total one.

\subsection{Race Conditions and Ordering}

Consider a classic race condition in a linked-list \texttt{insert} function. Two threads, $t_1$ and $t_2$, both call \texttt{insert} concurrently on a shared global list:

\begin{lstlisting}[language=C, caption={Linked list with race condition}]
struct list {
    int data;
    struct list *next;
};

struct list *list = 0;

void insert(int data) {
    struct list *l;
    l = malloc(sizeof *l);
    l->data = data;
    l->next = list;
    list = l;
}
\end{lstlisting}

Without synchronization, the interleaving of \texttt{l->next = list} and \texttt{list = l} between the two threads can produce a corrupted list in which one node is lost. The root cause is that the assignment to the global \texttt{list} pointer is not atomic.

The three levels of ordering relevant here are:

\begin{itemize}
    \item \textbf{Unordered}: Operations from $t_1$ and $t_2$ may interleave in any order ($\text{Rt}_1, \text{Rt}_2, \text{Wt}_1, \text{Wt}_2$ in any sequence).
    \item \textbf{Partially ordered} (with locks): Locks enforce that within each critical section the operations are totally ordered, but the two critical sections may execute in either sequence ($\text{Rt}_1, \text{Wt}_1, \text{Rt}_2, \text{Wt}_2$ or $\text{Rt}_2, \text{Wt}_2, \text{Rt}_1, \text{Wt}_1$).
    \item \textbf{Fully ordered}: Exactly one globally visible sequence of all operations exists, as in sequential consistency.
\end{itemize}

Adding acquire/release around the critical section (the read and write of \texttt{list}) restores safety: only one thread modifies \texttt{list} at a time, preventing the lost-update race.

% ============================================================
\section{Threading: Definition and Motivation}
% ============================================================

\subsection{Why Threading?}

The most frequently cited reason for adopting threads is \emph{performance}: multiple threads can overlap computation and I/O, keeping CPUs busy. However, multi-processing can achieve similar throughput gains. Threading offers additional benefits that processes alone cannot:

\begin{itemize}
    \item \textbf{Parallelism with shared memory}: Threads within a process share the same address space, eliminating the overhead of inter-process communication for shared data. A secondary benefit is more efficient context switching compared to full process switches, because no page-table swap is required.
    \item \textbf{Concurrency on a single core}: Even on a uniprocessor, a multithreaded program makes it easier to organize and respond to asynchronous events. Multiple threads give the programmer the \emph{abstraction} of concurrency even when true hardware parallelism is absent.
\end{itemize}

Threading also introduces challenges: \textbf{synchronization} (preventing races and ensuring ordering) and the inherent difficulty of \textbf{testing and debugging} concurrent programs.

\subsection{Threading as an Abstraction of Concurrency}

At a high level, \textbf{threading is a mechanism}---specifically, a mechanism for concurrency. More precisely, threading is an \emph{abstraction} of concurrency that gives the programmer the ability to execute independent threads at the same time, whether or not the hardware actually runs them in parallel.

\subsection{Other Mechanisms of Concurrency}

Threading is not the only concurrency mechanism. Two alternatives are:

\begin{itemize}
    \item \textbf{Multiprocessing}: Multiple processes share compute resources (the CPUs) but \emph{not} memory. The key guaranteed property is \textbf{memory space isolation}---one process cannot accidentally (or maliciously) corrupt another's address space.
    \item \textbf{Event-based systems}: A single process with a single thread voluntarily yields control at well-defined points (e.g., before blocking I/O). Node.js is a well-known example. This model is simple and avoids data races, but it provides no true hardware parallelism.
\end{itemize}

\subsection{Threads as a Mechanism of Concurrency}

The key design insight is: \emph{what if we deliberately break the memory isolation guarantee that processes provide?} The answer is threads. The \textbf{key insight} is that threads \emph{share memory resources by design}. This shared address space is the defining property that distinguishes a thread from a process.

\subsection{Defining Threads}

Two complementary definitions from the literature:

\begin{enumerate}
    \item A \emph{thread of execution} is the smallest sequence of programmed instructions that can be managed independently by a scheduler.
    \item A thread, sometimes called a \emph{lightweight process}, is a basic unit of CPU utilization; it comprises a thread ID, a program counter, a register set, and a stack.
\end{enumerate}

The defining structural property is: \textbf{threads within the same process share the same memory (address space, heap, globals)}, while each thread has its own stack and register state.

\subsection{Benefits of Threads}

Four classical benefits of threading are:

\begin{itemize}
    \item \textbf{Responsiveness}: Threads are especially beneficial for interactive applications. A foreground thread can respond to user input while background threads perform long-running work, providing the abstraction of concurrency to the user.
    \item \textbf{Resource sharing}: Threads share memory and resources by design. In contrast, processes have separate address spaces; sharing data between them requires explicit IPC (pipes, shared memory segments, etc.).
    \item \textbf{Economy}: Thread creation and destruction is much faster than process creation. Context switching between threads has less overhead because there is no page-table switch.
    \item \textbf{Multiprocessors}: On a system with multiple CPUs, threads can achieve \emph{true} parallelism rather than merely the single-processor illusion of time-sharing.
    \item Note, it's not that we can eliminate context switching, we just have less overhead.
\end{itemize}

\subsection{Design Considerations for Supporting Threads}

Providing threading support requires answering several questions:
\begin{itemize}
    \item How do we best structure and enable threading in the operating system?
    \item What is the interface for thread creation, destruction, and management?
    \item What is the real difference between a thread and a process? Does there \emph{need} to be one?
\end{itemize}

The last question is particularly interesting from an OS design perspective. Linux, for instance, treats threads and processes as the same underlying abstraction (\emph{tasks}), differing only in the degree of resource sharing.

% ============================================================
\section{Threading Design Tradeoffs}
% ============================================================

The central design question is how \textbf{user-level threads} (the programming abstraction) map onto \textbf{kernel-level threads} (the schedulable entities the OS manages). Three classical models exist.

\subsection{One-to-One (1:1) --- Kernel Threading}

In the \textbf{one-to-one} model, each user thread is backed by exactly one kernel thread. This is also called \emph{kernel threading} and is the model used by Linux (via pthreads/NPTL), Windows, and most modern operating systems.

\textbf{Attributes:}
\begin{itemize}
    \item \emph{Improved concurrency}: Multiple threads can run truly in parallel on a multiprocessor.
    \item \emph{Can run in parallel}: The OS scheduler can assign different kernel threads to different CPUs simultaneously.
    \item \emph{Additional overhead}: Creating a user thread requires creating a kernel thread, which involves a system call and kernel data structures.
\end{itemize}

\subsection{Many-to-One --- User Space Threading}

In the \textbf{many-to-one} model, many user-level threads are multiplexed onto a single kernel thread by a user-space thread library. The kernel is completely unaware of the individual user threads and schedules the entire process as a single entity.

\textbf{Attributes:}
\begin{itemize}
    \item Maps many user threads to one kernel thread; simple and efficient at the library level.
    \item True concurrency is difficult: if one user thread makes a blocking system call, it blocks the \emph{entire} kernel thread, freezing all other user threads.
    \item Only one thread can access the kernel at a time; no parallelism on multiprocessors.
\end{itemize}

\subsubsection*{How a User Space Threading Library Works}

A user-space thread library works roughly the same way xv6 manages kernel threads internally: it uses a \texttt{swtch()}-style context switch to swap register state, a scheduler to choose which user thread runs next, and concurrency primitives built from atomics and queues---all without any hardware protection boundaries. Crucially, such a library can still implement most of the threading machinery in user space. The question is what primitives it \emph{cannot} implement without kernel help.

The blocking-I/O problem illustrates the limit. If Thread~1 calls \texttt{sleep(100)} and Thread~2 calls \texttt{read(network, buf)}, a many-to-one library has no way to run Thread~1's remaining work while Thread~2's \texttt{read} system call is blocking inside the kernel. The kernel sees only one process and blocks it entirely.

\subsubsection*{User Space Threading Tradeoffs}

\begin{center}
\begin{tabular}{ll}
\toprule
\textbf{Pros} & \textbf{Cons} \\
\midrule
More efficient: no heavy kernel constructs & Cannot use 2+ CPUs without kernel help \\
No user-kernel switch for thread data & Needs kernel support for preemptive threading \\
More flexible per-application scheduling & No direct control over sleeping \\
 & File I/O can freeze all threads \\
\bottomrule
\end{tabular}
\end{center}

\subsection{Kernel Space vs.\ User Space Threading}

The choice between where to manage threads involves a direct tradeoff of capabilities:

\begin{center}
\begin{tabular}{ll}
\toprule
\textbf{Kernel Space} & \textbf{User Space} \\
\midrule
True parallelism & Semantically-aware (application-level) scheduling \\
Preemption possible & Lower context-switching overhead \\
More information available to scheduler & Better cache locality \\
Ability to control resource utilization & Lighter weight than kernel threading \\
 & Lower overhead of thread creation \\
\bottomrule
\end{tabular}
\end{center}

\subsection{Many-to-Many (M:N) --- Hybrid Threading}

The \textbf{many-to-many} model (also called \emph{M:N} or \emph{hybrid threading}) attempts to get the best of both worlds: $M$ user threads are multiplexed over $N$ kernel threads, where $N \leq M$. The user-space library schedules user threads onto available kernel threads, and the OS schedules the kernel threads onto CPUs.

\textbf{Advantages of M:N:}
\begin{itemize}
    \item User-space threads can run \emph{in parallel} (unlike pure many-to-one).
    \item Fewer kernel context switches than one-to-one.
    \item Lightweight user-space thread creation (small per-thread cost).
\end{itemize}

\textbf{Disadvantages:}
\begin{itemize}
    \item Still subject to kernel thread stalls: a blocking system call on one kernel thread can block all user threads mapped to it.
    \item Still not fully cooperative use of the CPU---the user-space scheduler has limited visibility into kernel scheduling decisions.
\end{itemize}

\subsection{Modern Threading Libraries}

Most modern threading libraries (e.g., POSIX pthreads on Linux) use \emph{kernel} threading (1:1), but implement many operations in user space to reduce overhead. The key kernel primitive enabling this is the \textbf{\texttt{futex}} (Fast Userspace muTEX) system call:

\begin{itemize}
    \item A futex is essentially a kernel sleep/wakeup mechanism exposed to user space.
    \item It allows user space to implement locking---checking and setting an atomic integer---entirely in user space when there is no contention. Only when contention occurs does the thread call into the kernel to sleep or wake another thread.
    \item This creates an efficient tradeoff: the \emph{kernel} manages I/O sleeping, process sleeping, and scheduling, while \emph{user space} manages most concurrency primitives (locks, condition variables, semaphores).
\end{itemize}

\subsection{Scheduler Activations}

\textbf{Scheduler activations} (Anderson et al., 1991) represent one concrete approach to realizing M:N scheduling. The core idea is to introduce a \emph{communication mechanism} between kernel space and user space:

\begin{itemize}
    \item The kernel notifies the application's user-space thread scheduler of kernel events that affect the address space (e.g., a thread blocking on I/O, a new CPU being assigned). This notification is called an \emph{upcall}.
    \item In response, the user-space scheduler can modify its thread data structures, run another thread, and make requests back to the kernel.
\end{itemize}

Scheduler activations were implemented in NetBSD but ultimately abandoned, and were also attempted in Linux 2.4.x. The complexity of the kernel--user communication protocol---and the difficulty of handling all edge cases correctly---made the approach difficult to maintain in practice.

\subsubsection*{Modern Renditions of Scheduler Activations}

The idea has seen a resurgence in systems research:
\begin{itemize}
    \item \textbf{ghOSt} (SOSP'21, Google): A framework in which the kernel delegates CPU scheduling decisions for a set of threads to a user-space agent process. The kernel sends thread/CPU messages to the agent via shared memory; the agent makes scheduling decisions and sends them back via transactions and system calls.
    \item \textbf{Arachne} (OSDI'18, Stanford): A user-level threading system in which a core arbiter communicates with each application using shared memory and sockets. Applications can use custom thread libraries and per-application scheduling policies layered on top of the arbiter.
\end{itemize}

Both systems illustrate how the semantic gap between kernel and application can be addressed without fully committing to user-space threading.

\subsection{Why Kernel Threads Have (Mostly) Won}

Despite the theoretical appeal of user-space and hybrid threading, the 1:1 kernel-threading model has become the near-universal standard. Three reasons explain this:

\begin{enumerate}
    \item \textbf{Blocking system calls}: Multiple user threads sharing a single kernel thread are inherently limited. A single blocking system call (e.g., a network \texttt{read}) blocks the entire kernel thread and therefore all user threads mapped to it.
    \item \textbf{Semantic gap / poor scheduling decisions}: The kernel has no awareness of user-level thread structure, leading to suboptimal scheduling. This is sometimes called the \emph{semantic gap} or lack of \emph{semantic awareness}.
    \item \textbf{True parallelism}: With pure user-space threading, each process has access to at most one CPU at a time, making it impossible to exploit multiple cores within a single process.
\end{enumerate}

% ============================================================
\section{Threading Support in Modern Operating Systems}
% ============================================================

\subsection{Threading Models Across Systems}

The evolution of threading support varies substantially across operating systems. Most modern systems have converged on the 1:1 kernel-threading model:

\begin{center}
\begin{tabular}{llll}
\toprule
\textbf{Operating System} & \textbf{Introduced} & \textbf{Model} & \textbf{Notes} \\
\midrule
Minix          & 1987  & N/A           & It's complicated \\
Linux (LinuxThreads) & 1996 & 1:1      & \\
Linux (NPTL)   & 2003  & 1:1           & \\
macOS/iOS      & $\sim$2000 & 1:1ish   & Built on Mach Threads, BSD layer \\
FreeBSD        & 2003  & M:N or 1:1    & 1:1 default/recommended since 2006 \\
DragonFly BSD  & 2007  & 1:1           & \\
NetBSD         & 2009  & 1:1           & 1:1 since 5.0 (circa 2009) \\
OpenBSD        & 2012  & 1:1           & Replaced M:N \\
Solaris        & 2000  & M:N then 1:1  & \\
HP-UX          & 1997  & M:N or 1:1    & \\
Windows        & 1993  & 1:1           & \\
Fuchsia        & 2016  & 1:1           & \\
\bottomrule
\end{tabular}
\end{center}

The historical trend is clear: systems that initially adopted M:N threading (NetBSD, OpenBSD, Solaris, FreeBSD) have largely migrated to 1:1, confirming that the complexity of M:N rarely justifies the theoretical benefits in practice.

\subsection{Threads in Linux: A Concrete Example}

The following example creates \texttt{NUM\_THREADS} POSIX threads using \texttt{pthread\_create}, each of which prints its process ID and thread ID before sleeping:

\begin{lstlisting}[language=C, caption={Hello World with pthreads (adapted from Purdue ECE563)}]
#define NUM_THREADS 5

void *TaskCode(void *argument) {
    int tnum;
    pid_t pid = getpid();
    pid_t tid = syscall(SYS_gettid);

    tnum = *((int *) argument);
    printf("Hello World from thread #%d! My PID is %d and my TID is %d\n",
           tnum, pid, tid);
    sleep(600);
    return NULL;
}

int main(void) {
    pthread_t threads[NUM_THREADS];
    int thread_args[NUM_THREADS];
    int rc, i;
    for (i = 0; i < NUM_THREADS; ++i) {
        thread_args[i] = i;
        printf("[Parent PID %d]: In main, creating thread %d\n", getpid(), i);
        rc = pthread_create(&threads[i], NULL, TaskCode, (void *) &thread_args[i]);
        assert(0 == rc);
    }
}
\end{lstlisting}

Notably, all threads share the same PID (they belong to the same process), but each has a distinct TID (thread group ID). Running \texttt{ps -T -o pid,tgid,tid,cmd -p \textit{PID}} from another terminal shows all six kernel threads (the main thread plus five created threads) with the same PID and TGID but sequential TIDs---confirming the 1:1 mapping between user threads and kernel threads in Linux.

% ============================================================
\section{Threading in xv6}
% ============================================================

\subsection{xv6's Current Threading Model}

xv6 does not currently support user threads. Its threading model is minimal:
\begin{itemize}
    \item \textbf{Kernel stack}: Each process has its own kernel stack, which serves as the kernel-side execution context.
    \item \textbf{User threads}: There are none. Each xv6 process is a single-threaded entity from the user's perspective.
    \item \textbf{Scheduler awareness}: xv6's scheduler has no awareness of symmetric multiprocessing (SMP). It has no SMP support; the scheduler is effectively uniprocessor.
\end{itemize}

\subsection{Adding Kernel Thread Support to xv6}

To extend xv6 with 1:1 kernel thread support (as in Lab 3), three categories of additions are needed:

\begin{enumerate}
    \item \textbf{Schedulability}: Threads must be visible to the scheduler---meaning each thread must have its own kernel-side context (stack, register set) and appear in the process table as a runnable entity.
    \item \textbf{Performance advantages}: Threads must share memory. This means sharing the page directory (\texttt{pgdir}) so all threads within a process map to the same physical address space.
    \item \textbf{Synchronization primitives}: Threads that share memory need coordinated access. The required primitives are:
    \begin{itemize}
        \item \emph{Waiting}: \texttt{waitpid}, \texttt{thread\_wait}
        \item \emph{Spinlocks}: \texttt{init}, \texttt{acquire}, \texttt{release}
        \item \emph{Mutexes}: \texttt{park}/\texttt{setpark}/\texttt{unpark}, mutex \texttt{init}/\texttt{acquire}/\texttt{release}
        \item \emph{Condition variables}: \texttt{cond\_wait}, \texttt{cond\_signal}
    \end{itemize}
\end{enumerate}

\subsection{The \texttt{clone()} System Call}

The mechanism for creating a thread in xv6 is a new system call, \texttt{clone()}, modeled after the Linux \texttt{clone()} system call. In Linux, \texttt{clone()} is the single underlying primitive that implements \emph{both} thread and process creation:

\begin{itemize}
    \item \texttt{fork()} creates a new process with its \emph{own} memory space (copy-on-write). Internally, Linux implements \texttt{fork()} via \texttt{clone()} with flags that specify a new memory space.
    \item \texttt{clone()} with \texttt{CLONE\_VM} creates a new task that \emph{shares} the parent's memory space---this is a thread.
\end{itemize}

Every Linux task is a schedulable entity; whether it is a ``process'' or a ``thread'' is determined entirely by which resources it shares with its creator.

The Linux \texttt{clone()} system call signature is:
\begin{lstlisting}[language=C, caption={Linux clone() system call}]
long clone(unsigned long flags, void *child_stack,
           void *ptid, void *ctid,
           struct pt_regs *regs);
\end{lstlisting}

And the library-level wrapper exposed to user programs:
\begin{lstlisting}[language=C, caption={Linux clone() library function}]
int clone(int (*fn)(void *), void *child_stack,
          int flags, void *arg, ...
          /* pid_t *ptid, struct user_desc *tls, pid_t *ctid */ );
\end{lstlisting}

The extensive list of \texttt{CLONE\_*} flags (including \texttt{CLONE\_VM}, \texttt{CLONE\_FILES}, \texttt{CLONE\_SIGHAND}, \texttt{CLONE\_THREAD}, and many others) controls exactly which resources the new task shares with its parent. This flag-based design is what allows a single system call to express the full spectrum from a fully isolated new process to a fully shared new thread.

\subsection{Clone in xv6}

For Lab 3, xv6 uses a simplified \texttt{clone()} interface:

\begin{lstlisting}[language=C, caption={xv6 clone() interface}]
int clone(void *stack, int size);
\end{lstlisting}

The implementation is similar to \texttt{fork()}, with one critical difference: instead of duplicating the parent's address space (with copy-on-write), \texttt{clone()} in xv6 \emph{reuses the parent's page directory entirely}. Both the parent and the new thread thus share the same physical memory mapping. The \texttt{stack} argument allows the caller to supply a pre-allocated stack region for the new thread to use.

\subsection{Pthreads in xv6}

The standard POSIX thread creation interface is \texttt{pthread\_create}:

\begin{lstlisting}[language=C, caption={POSIX pthread\_create interface}]
int pthread_create(pthread_t *thread,
                   const pthread_attr_t *attr,
                   void *(*start_routine)(void *),
                   void *arg);
\end{lstlisting}

For xv6, a simpler \texttt{thread\_create} wrapper is provided that hides stack allocation and the \texttt{clone()} call:

\begin{lstlisting}[language=C, caption={xv6 thread\_create interface}]
thread_create(void *(*start_routine)(void *),
              void *arg)
\end{lstlisting}

This wrapper allocates a stack page, calls \texttt{clone(stack, size)}, and arranges for the new thread to begin execution at \texttt{start\_routine(arg)}, presenting the programmer with an interface analogous to POSIX pthreads.

% ============================================================
\section{Summary}
% ============================================================

\begin{itemize}
    \item \textbf{Weak ordering}: Hardware and compilers may reorder memory operations for performance; without explicit synchronization, cross-thread ordering guarantees are almost nonexistent. Lock acquire/release provides partial ordering.
    \item \textbf{Threads vs.\ processes}: The defining property of a thread is \emph{shared memory by design}. Threads share address space; processes are isolated. This tradeoff enables efficient communication at the cost of requiring synchronization.
    \item \textbf{Benefits of threading}: Responsiveness, resource sharing, economy of creation/switching, and potential true parallelism on multiprocessors.
    \item \textbf{One-to-One (1:1)}: Each user thread is backed by a kernel thread. Enables true parallelism and correct blocking semantics. Standard in Linux (NPTL), Windows, macOS, and most modern OSes.
    \item \textbf{Many-to-One}: User threads multiplexed on one kernel thread. Efficient but no parallelism and blocking calls freeze all threads.
    \item \textbf{Many-to-Many (M:N)}: Hybrid model offering partial parallelism and efficiency. Complex to implement correctly; most systems have abandoned it in favor of 1:1.
    \item \textbf{Scheduler activations}: A communication mechanism between kernel and user-space scheduler. Theoretically elegant; abandoned in practice due to complexity. Modern relatives include ghOSt and Arachne.
    \item \textbf{Kernel threads have mostly won}: Blocking syscalls, the semantic gap, and the need for true parallelism make 1:1 the practical winner.
    \item \textbf{\texttt{clone()} in Linux and xv6}: Both threads and processes are created via \texttt{clone()}, which unifies them as tasks differing only in their degree of resource sharing. The xv6 Lab 3 \texttt{clone()} reuses the parent's page directory to achieve shared memory.
\end{itemize}

% ============================================================
\section{References}
% ============================================================

\begin{itemize}
    \item Anderson, T. E. et al. ``Scheduler Activations: Effective Kernel Support for the User-Level Management of Parallelism.'' \textit{ACM Transactions on Computer Systems}. 1992. \url{https://web.eecs.umich.edu/~mosharaf/Readings/Scheduler-Activations.pdf}
    \item Silberschatz, A.; Galvin, P. B.; Gagne, G. (1999). \textit{Applied Operating System Concepts}. John Wiley \& Sons.
    \item Ghavami, B. et al. ``ghOSt: Fast \& Flexible User-Space Delegation of Linux Scheduling.'' \textit{SOSP 2021}.
    \item Qin, H. et al. ``Arachne: Core-Aware Thread Management.'' \textit{OSDI 2018}.
    \item Wikipedia contributors. ``Thread (computing).'' \textit{Wikipedia}. \url{https://en.wikipedia.org/wiki/Thread_(computing)}
    \item Adapted pthreads example from: \url{https://engineering.purdue.edu/~eigenman/ECE563/Handouts/helloworld-pthreads.c}
    \item Unix history diagram: Eraserhead1, Infinity0, Sav\_vas. Levenez Unix History Diagram. CC BY-SA 3.0. \url{https://commons.wikimedia.org/w/index.php?curid=1801948}
\end{itemize}

% ============================================================
\section*{Personal Notes}
% ============================================================
% Add your own notes below this line

\end{document}
